\section{Research objective and scope}
\label{sec:final-research-objective}

This thesis investigates the detection capability and generalization
boundaries of interpretable audio phenomena in AI-generated and human music.
Its objective is not to infer authorship from a single presumed acoustic
signature. Instead, it asks whether explicitly defined measurements respond
to their intended phenomena, provide incremental discrimination within a
documented corpus, and retain that information when recording sources or
generators are held out.

The motivating observation is perceptual: high-frequency characteristics,
particularly in vocals, can influence a listener's impression that music was
generated. Turning this observation into a scientific claim requires more
than a classifier trained on two collections. Frequency response, encoding,
mastering, instrumentation, vocal activity and source separation can change
the same measurements. A feature can be interpretable in its mathematical
definition without uniquely identifying the physical process that produced
an observed value.

\subsection{Three distinct evidential questions}
\label{sec:three-evidence-levels}

The study separates three questions. First, \emph{measurement validity}:
does a descriptor detect a controlled change in its intended phenomenon,
and with what availability and nuisance sensitivity? Second,
\emph{discriminative utility}: does adding the descriptor improve a fixed
classification procedure on matched recordings and evaluation splits?
Third, \emph{generalization}: does the resulting discrimination persist
outside the recording sources and generators represented in training?
These questions are related but not interchangeable. A passed external
measurement gate does not establish classification utility; a high
classification score does not retrospectively validate a failed measurement
gate; and ordinary grouped evaluation does not establish unseen-generator
attribution.

\subsection{Operational meaning of the labels}

Human and AI labels refer to the documented provenance of selected recordings.
They do not label every musical decision or component in a production chain.
Human-produced music may contain MIDI programming, synthesis, quantization,
editing and compression. Consequently, regular dynamics or rhythm cannot
alone establish AI generation. Conversely, a generated recording can exhibit
variation in these quantities. The proposed measurements are statistical
descriptors of audio, not necessary or sufficient conditions for either class.
Provisional labels and small conditioning-group pilots are kept separate from
validated accuracy claims.

\subsection{Research questions}

\begin{enumerate}
\item Which listening-motivated spectral differences remain observable after
      standardized audio preparation and an external separator-response
      correction, and what confounds remain unresolved?
\item Which additional phenomena can be measured reliably under the retained
      external controlled tests, including explicit negative controls and
      missingness accounting?
\item How do individual feature families and their nonempty combinations
      compare under fixed, source-aware splits and training-group caps?
\item Where do source-held-out and generator-held-out results diverge from
      ordinary grouped performance, and what does this imply for the scope
      of the detector's claims?
\end{enumerate}

\subsection{Reporting commitments}

All conclusions are tied to named cohorts, preprocessing versions and
evaluation protocols. Recording identities, independent groups, audio views
and stored files are different units and are reported separately. Comparisons
across different duration cohorts are descriptive rather than causal estimates
of duration effects. Missing measurements remain explicit; absence of a
measurement is not silently interpreted as absence of the phenomenon.

The thesis retains unsuccessful candidates and negative feature increments.
Numerical audits establish only the properties that they explicitly check;
byte-level integrity is distinguished from waveform replay, and both are
distinguished from external scientific validity. Raw classifier scores are
not presented as calibrated probabilities. Existing development sweeps are
not relabelled as untouched validation after inspecting their rankings.

The final experimental scope is restricted to the already collected material
and the existing remaining analysis plan. New corpora, new phenomena and
additional experimental branches are outside this thesis finalization run.
Unresolved weaknesses are discussed as limitations and future work rather
than addressed by unreported post-hoc experiments.
